% ==========================================================
\section{ANTLR Parser Integration and Metamodel Extensions for BP}
\label{sec:impl_bp_java}
% ==========================================================

This section describes the concrete realization of the conceptual BP integration introduced in Section~\ref{sec:concept_bp_integration}, covering the grammar extensions, parser adaptations, and metamodel classes that constitute the implementation.
Whereas Section~\ref{sec:concept_bp_integration} motivates the design decisions, this section focuses on how those decisions are realized in the code base.
The three implementation layers form a directed pipeline: the grammar extensions define what syntax is legal, the parser listener translates parsed input into objects, and the metamodel classes provide the structure those objects conform to.

\subsection{ANTLR Grammar Extensions}
\label{subsec:antlr_extensions}

The grammar modifications are confined to \texttt{UVLJavaParser.g4} in the ANTLR grammar module of the UVL metamodel parser component.
They realize the opt-in BP extension described at the concept level while leaving all standard UVL productions unchanged, thereby preserving backward compatibility with non-BP models.
Because the UVL metamodel repository is already built on ANTLR-generated parsing, extending the grammar is the natural mechanism for introducing new syntax.
This stands in contrast to the UVLS implementation by Felber~\cite{felber_comprehensible_2025}, which is based on tree-sitter and required no grammar modifications to support BP constructs.

The extensions comprise three targeted additions, shown in Listing~\ref{lst:antlr-extension}.
First, the \texttt{features} rule is augmented to allow up to two additional top-level features, \texttt{Env} and \texttt{Config}, which represent the environmental and configuration contexts required by behavioral programs.
Second, the \texttt{constraint} rule is extended with a new labeled alternative \texttt{bpConstraint}, which groups all BP-specific constraint types under a single production to maintain a clean grammar structure.
Third, within \texttt{bpConstraint}, a dedicated \texttt{conflictingConstraint} production handles the multi-reference case of mutually conflicting events.
Listing~\ref{lst:antlr-extension} shows the relevant grammar excerpt.

\lstinputlisting[
    float,
    floatplacement=ht,
    caption={[ANTLR grammar extensions]ANTLR grammar extensions for BP integration},
    label={lst:antlr-extension}
]{code/metamodel/uvl_bp_grammar_snippet.g4}

To ensure deterministic lexing, each BP construct is introduced with a dedicated keyword token.  Listing~\ref{lst:bp-lexer-extension} shows the token definitions.

\lstinputlisting[
    float,
    floatplacement=ht,
    caption={[ANTLR lexer extensions]ANTLR lexer extensions for BP integration},
    label={lst:bp-lexer-extension}
]{code/metamodel/bp_lexer_extensions.g4}

\subsection{Parser and Factory Adaptations}
\label{subsec:parser_factory_edits}

The parsing logic is adapted in two places: the ANTLR listener and the model factory.
Model variant selection is controlled by the \texttt{ModelType} enum, which was introduced to support distinct parsing strategies under a single parser entry point and may be extended to accommodate additional model variants in the future.
When \texttt{ModelType.BP} is active, the \texttt{UVLListener} constructor instantiates a \texttt{BPFeatureModel}.
Otherwise, it falls back to the standard \texttt{FeatureModel}, as shown in Listing~\ref{lst:uvl-listener-constructor}.

\lstinputlisting[
    language=Java,
    float,
    floatplacement=ht,
    caption={[\texttt{UVLListener} constructor]\texttt{UVLListener} constructor with model type selection.},
    label={lst:uvl-listener-constructor}
]{code/metamodel/UVLListener_constructor.java}

ANTLR generates paired \texttt{enterX} and \texttt{exitX} callbacks for each grammar rule and labeled alternative.
A \texttt{ParseTreeWalker} invokes \texttt{enterX} before visiting a node's children and \texttt{exitX} after all child callbacks have completed.
This separation is exploited systematically: \texttt{enter} methods initialize objects and push them onto a stack, while \texttt{exit} methods finalize those objects using data populated by child callbacks~\cite{parr_definitive_2013}.

The BP top-level feature rule illustrates this pattern concretely.
Its grammar production is \texttt{reference attributes? NEWLINE}, meaning that when the walker enters the rule, the \texttt{reference} and \texttt{attributes} subtrees are accessible but not yet resolved.
Both listener methods begin with a call to \texttt{rejectIfBaseModel}, which enforces the opt-in nature of BP.
When the active model type is not \texttt{ModelType.BP}, any BP-specific construct encountered in the listener is rejected with a parse error and the method returns immediately.
This guard is necessary because the grammar is shared across model types.
The distinction between permitted and forbidden constructs is therefore enforced at the listener level rather than in the grammar itself~\cite{TODO_sourcecode}.

Accordingly, \texttt{enterEnvConfigFeature} creates a \texttt{Feature} and pushes it onto the \texttt{featureStack}.
Child listeners then populate the new feature with any nested data, which in this case is just attributes.
When \texttt{exitEnvConfigFeature} is invoked, the feature is popped from the stack and assigned to either the \texttt{Env} or \texttt{Config} field in the \texttt{BPFeatureModel}, depending on its name.
Any feature name other than these two is rejected with a parse error.
The corresponding listener methods are shown in Listing~\ref{lst:bp-top-level-listener}.

\lstinputlisting[
    language=Java,
    float,
    floatplacement=ht,
    caption={[BP top-level feature listener]Listener methods for BP top-level features (Env / Config).},
    label={lst:bp-top-level-listener}
]{code/metamodel/BP_top_level_listener.java}

BP constraints follow a similar enter/exit pattern.
The enter method instantiates the appropriate constraint object, and the exit method finalizes it and registers it with the model.
For \texttt{ConflictingConstraint}, which carries multiple references, the listener accumulates all reference objects into a list during child callbacks and passes that list to the constructor upon exit.
The full implementations are available in the accompanying source code~\cite{TODO_sourcecode}.

\subsection{Metamodel Extensions}
\label{subsec:metamodel_extensions}

The metamodel classes introduced in this subsection are the same ones instantiated by the listener described in Section~\ref{subsec:parser_factory_edits}.
They provide the structural definitions that the listener populates during parsing.

The central addition is \texttt{BPFeatureModel}, which extends \texttt{FeatureModel} and introduces two dedicated fields for the \texttt{Env} and \texttt{Config} top-level features.
Both are represented as standard \texttt{Feature} objects and are distinguished semantically by their role within the BP context rather than by a separate type.
Listing~\ref{lst:bp-feature-model} shows the class structure, including the override of \texttt{appendAdditionalTopLevelFeatures}, which is called by the base class \texttt{toString} method during serialization and ensures that \texttt{Env} and \texttt{Config} are included in the serialized model output.

\lstinputlisting[
    language=Java,
    float,
    floatplacement=ht,
    caption={[\texttt{BPFeatureModel} class]\texttt{BPFeatureModel} class with \texttt{Env} and \texttt{Config} top-level features.},
    label={lst:bp-feature-model}
]{code/metamodel/BPFeatureModel.java}

Another example of a metamodel addition is \texttt{ConflictingConstraint}, which extends \texttt{Constraint} to represent conflicts between multiple events.
Rather than holding a single reference, it manages an ordered list of \texttt{VariableReference} objects.
The list is exposed as an unmodifiable view to prevent unintended mutation, while individual references can be replaced via an indexed setter to support reference updates after construction.
The class is shown in Listing~\ref{lst:conflicting-constraint}.

\lstinputlisting[
    language=Java,
    float,
    floatplacement=ht,
    caption={[\texttt{ConflictingConstraint} class]\texttt{ConflictingConstraint} class for multi-reference BP constraints (excerpt).},
    label={lst:conflicting-constraint}
]{code/metamodel/ConflictingConstraint.java}

All remaining BP constraint types follow the same structural pattern: a grammar production, a listener that constructs the constraint object, and registration with the model upon exit.
Single-reference BP constraints extend \texttt{AbstractBPEventConstraint}, which is a shared base class that encapsulates the common single-event reference structure.
Further implementation details are available in the source code~\cite{TODO_sourcecode}.

In summary, the BP integration is realized through targeted, additive extensions at each layer of the pipeline.
The grammar extensions are the natural consequence of the UVL metamodel repository's existing dependency on ANTLR-generated parsing; the same integration achieved without grammar changes in the tree-sitter-based UVLS implementation by Felber~\cite{felber_comprehensible_2025} reflects an infrastructure difference rather than a fundamental distinction in approach.
All changes are strictly additive with respect to the standard UVL pipeline, and the result is a BP-capable parsing infrastructure that is backward compatible with existing UVL models and structurally aligned with the conceptual architecture developed in this thesis.